% Drawing 70/03-4 -- folder 70, batch 03, archive page 58 (PDF page 59).
% Transcription: un sommet $e$ d'où partent une arête vers un autre sommet et une flèche $f$ ; à côté, l'image de la situation en $\gamma e$.
% Kind: vertex with an edge and an arrow. Method 2 (schematic redraw from the facsimile). Status: ok.
% Notes: Vertex e with an edge (a curve) up to a marked vertex and a thick arrow f; to the right, gamma e with a long thin arc, the image of the edge. The horizontal rule under the figure separates it from the text and is not drawn. Redrawn 2026-11-09 (issue 128) after Michel's check: vertices placed once and every edge drawn between them, so concurrent lines meet exactly; labels at \small.
% Coordinates: 1 unit = 1 mm = 1/100 of the crop width (bbox in page fractions: [0.5, 0.24, 0.82, 0.36]).
% Redrawn by Claude (claude-opus-5-5), issue #73, 2026-11-09; not checked by a human; not compiled (no TeX in the session).
% Rights: reproduces part of the work -- behind RENDER_TRANSCRIPTS, never published by itself.
\documentclass[tikz,border=3pt]{standalone}
\usepackage{amsmath,amssymb}
\usetikzlibrary{patterns,arrows.meta}
\tikzset{pencil/.style={black!45}, struck/.style={black!35}, mark/.style={circle,draw,inner sep=1.3pt}, pt/.style={circle,fill,inner sep=0.9pt}, opt/.style={circle,draw,fill=white,inner sep=1.1pt}}
\begin{document}
\begin{tikzpicture}[x=1mm,y=1mm,line cap=round,line join=round,font=\scriptsize]
\draw plot[smooth, tension=0.55] coordinates {(10,-42.5) (12,-25) (18,-12.5)};
\node[pt] at (18,-12.5) {};
\node[pt] at (10,-42.5) {};
\draw[-{Stealth[length=2mm]},very thick] plot[smooth, tension=0.55] coordinates {(10,-42.5) (25,-33) (39,-24.5)};
\node[inner sep=1pt] at (11,-47.5) {\small $e$};
\node[inner sep=1pt] at (22,-30.5) {\small $f$};
\node[pt] at (60,-41) {};
\draw plot[smooth, tension=0.55] coordinates {(60,-41) (75,-29) (93,-21)};
\node[inner sep=1pt] at (65,-46) {\small $\gamma e$};
\end{tikzpicture}
\end{document}
